\documentclass[11pt,a4paper]{article}
\usepackage{fontspec,amsmath,amssymb,amsthm,microtype}
\setmainfont{lmroman10-regular.otf}[BoldFont=lmroman10-bold.otf,ItalicFont=lmroman10-italic.otf,BoldItalicFont=lmroman10-bolditalic.otf]
\setmonofont{lmmono10-regular.otf}[Scale=MatchLowercase]
\usepackage[margin=25mm]{geometry}
\usepackage[unicode,hidelinks]{hyperref}
\setlength{\emergencystretch}{3em}
\setlength{\parskip}{0.4em}
\setlength{\parindent}{0pt}
\allowdisplaybreaks[2]
\raggedbottom
\title{\texorpdfstring{The exact physical-\allowbreak{}time integral equation}{The exact physical-time integral equation}}
\author{LangevinLean mathematical exposition}
\date{}
\begin{document}\maketitle
\noindent Source module: \texttt{KIPLMC/\allowbreak{}Kinetic\allowbreak{}Clock.\allowbreak{}lean}. The accompanying HTML pairs this mathematical exposition with the exact Lean statements and records the scope of the formalization.


Let \(Y_s\) solve the unit-temperature kinetic equation with sum force \(g\) and friction \(c\gamma\). The physical process is \(S_t=D_{c^{-1}}Y_{t/c}\), where \(D_a(v,x)=(av,x)\). This module proves the complete transformed integral equation in any complete normed vector space, including the L\ensuremath{{}^{2}} space of actual random variables used by the Brownian construction.


\section{\texorpdfstring{Make velocity scaling a bounded linear map}{Make velocity scaling a bounded linear map}}


The velocity projection multiplied by \(c\), paired with the unchanged position projection, defines a continuous linear map. Its action is exactly \texttt{rescale\allowbreak{}Velocity}. Linearity allows this map to pass through Bochner integrals; continuity ensures the transformed path remains continuous. Expanding the phase drift proves \(c^{-1}D_{c^{-1}}R_{g,c\gamma}(z)=R_{c^{-2}g,\gamma}(D_{c^{-1}}z)\). The two powers of \(c^{-1}\) in the force come from time and velocity separately.



{\raggedright\noindent Lean declarations: \texttt{velocity\allowbreak{}Scale\allowbreak{}Linear}, \texttt{velocity\allowbreak{}Scale\allowbreak{}Linear\_\allowbreak{}apply}, \texttt{kinetic\allowbreak{}Negative\allowbreak{}Drift\_\allowbreak{}clock}.\par}

\section{\texorpdfstring{Change variables inside the actual integral equation}{Change variables inside the actual integral equation}}


For \(0\le t\le T\), positivity of \(c\) gives \(t/c\in[0,T/c]\). Apply the continuous linear map to the given solution identity at \(t/c\), pass it through the integral, and substitute \(r=s/c\). The drift identity identifies the transformed force and friction, and the additive source becomes \(c^{-1}B_{t/c}\). Lean proves equality of the complete vector-valued integrals, not merely equality of formal differential symbols. This theorem requires no derivative of the Brownian input.



{\raggedright\noindent Lean declarations: \texttt{kinetic\_\allowbreak{}integral\_\allowbreak{}clock}.\par}

\end{document}
